\chapter*{Kurzzusammenfassung}

   Bewegungsanalysesysteme erreichen zunehmend größere Bedeutung, zum Beispiel in der Medizintechnik in Diagnoseapplikationen oder als Mensch-Maschine-Schnittstelle in VR/AR-Reality-Anwendungen. Dabei sind besonders Sensorhandschuhe, mit denen die Bewegung einer Hand aufgenommen werden können von besonderem Interesse. Der Anstoß für diese Arbeit war es, einen Sensorhandschuh für die Steuerung eines roboterassistierten Chirurgiesystem zu verwenden.
   Da der menschliche Körper magnetisch neutral ist und die Felder von diesem nicht verzerrt werden, werden magnetische Sensorhandschuhe erforscht. Im Vergleich zu optischen Systemen benötigen diese keinen direkten Sichtkontakt, was bei komplexen kinematischen Ketten teilweise schwierig ist. 
   In dieser Arbeit wurde ein System zur Handbewegungsanalyse entwickelt. Dieses basiert auf zwei unabhängigen Lokalisierungsmethoden, eine wird verwendet um die Pose der Finger im Bezug zum Handrücken geschätzt und eine um die absolute Position der Hand im Raum zu bestimmen. Für beide Lokalisierungen werden 3D-Spulen verwendet die einen rotierenden magnetischen Dipol erzeugen. Es wird ein neues Signalmerkmal eingeführt, der Maximalvektor. Dieser definiert die Richtung in der an einem Sensor das maximale Feld detektiert wird. Der Maximalvektor wird über eine räumliche Orthogonalität bestimmt in dem zwei Nulldurchgänge im Signal detektiert werden. 
   Für die Schätzung der kinematischen Kette wurde ein Algorithmus entwickelt, der Wissen über die menschliche Anatomie und die Ausbreitung magnetischer Felder miteinander verknüpft. So kann über ein iteratives Verfahren effizient die Handhaltung bestimmt werden.
   Für die absolute Lokalisierung wurde ein modularer Algorithmus entwickelt, mit dem zunächst die drei Positionskomponenten der 3D-Spule einzeln geschätzt werden und darauf aufbauend die Orientierung. 
   Die Schätzung der kintematischen Kette wurde experimentell mit einem einfachen Prototyp bestehend aus drei kinematischen Elementen. Die statischen Schätzungen haben einen mittleren Fehler von $1.3^\circ/2.7^\circ$. In den dynamischen Messungen wurde ein Fehler von $1.9^\circ/5.6^\circ$ erreicht.
   Die absolute Lokalisierung hat in der experimentelle Positionsschätzung im Mittel einen Fehler von 1.5\,cm im statischen Fall und 1.7\,cm dynamischen Fall. Die Orientierungsschätzung hat einen mittleren Fehler von $15^\circ$ im statischen und $26^\circ$ im dynamischen Fall. 
   Mit der Evaluierung konnte qualitativ gezeigt werden, dass die Bewegungsanalyse möglich ist.
   
   

\chapter*{Abstract}

Motion analysis systems are becoming increasingly important, for example in medical technology in diagnostic applications or as human-machine interfaces in VR/AR reality applications. Sensor gloves that can be used to record the movement of a hand are of particular interest in this context. The impetus for this work was to use a sensor glove to control a robot-assisted surgical system.
Since the human body is magnetically neutral and does not distort magnetic fields, magnetic sensor gloves are being researched. Compared to optical systems, these do not require direct visual contact, which can be difficult in some cases with complex kinematic chains. 
In this work, a system for hand movement analysis was developed. It is based on two independent localisation methods, one used to estimate the pose of the fingers in relation to the back of the hand and one to determine the absolute position of the hand in space. For both localisations, 3D coils are used to generate a rotating magnetic dipole. A new signal feature is introduced, the maximum vector. This defines the direction in which the maximum field is detected at a sensor. The maximum vector is determined via spatial orthogonality in which two zero crossings are detected in the signal. An algorithm was developed to estimate the kinematic chain, combining knowledge of human anatomy and the propagation of magnetic fields. This allows the hand position to be determined efficiently using an iterative process.
For absolute localisation, a modular algorithm was developed that first estimates the three position components of the 3D coil individually and then uses this information to estimate the orientation. 
The estimation of the kinematic chain was tested experimentally with a simple prototype consisting of three kinematic elements. The static estimates have a mean error of $1.3^\circ/2.7^\circ$. In the dynamic measurements, an error of $1.9^\circ/5.6^\circ$ was achieved. Absolute localisation has an average error of 1.5 cm in static conditions and 1.7 cm in dynamic conditions in experimental position estimation. Orientation estimation has an average error of $15^\circ$ in static conditions and $26^\circ$ in dynamic conditions. The evaluation qualitatively demonstrated that motion analysis is possible.

